\section{Сертифицированная замена контроллера}\label{sec:replacement}
Предыдущие конструкции отделяют существование переносимого исследования от стоимости его обнаружения. Однако выбор контроллера до начала поиска не обосновывает его замену в ходе частично завершённого поиска. Замена должна сохранять и неразрешённые обязательства, и уже сертифицированные результаты; её оплата не должна лишать возможности завершить исследование. Зафиксируем индексированные исходную и целевую задачи $S,T$ и канонический эталонный список $\CanonRef=\rc{discoverAll}(S,T)$ раздела~3. В разделах~6--7 учётные величины означают натуральные модельные платы, а не измеренное время или память.

\subsection{Продолжения и их денотация}
Контроллер $C=(\mathcal R,\tau)$ состоит из библиотеки сертифицированных правил отсечения и базовых тарифов поиска. Фронтир $f=(J,A,\Lambda)$ содержит ожидающие задания $J$, список сертифицированных морфизмов $A$ и список поисковых квитанций $\Lambda$. Обозначая через $\Cov(x,J)$ покрытие морфизма ожидающим заданием, определим корректность фронтира в точности как
\begin{equation}\label{eq:frontier-correct}
 \Correct(f)\;\Longleftrightarrow\;
 (A\subseteq\CanonRef)\ \land\
 \bigl(\forall x\in\CanonRef,\ x\in A\ \lor\ \Cov(x,J)\bigr),
\end{equation}
где включение и членство относятся к спискам и не утверждают уникальность или непересечение. Таким образом, фронтир представляет и установленное знание, и оставшиеся обязательства поиска. Для фронтиров $f,g$ определим эквивалентность продолжений:
\begin{equation}\label{eq:same-cont}
 f\simeq g\;\Longleftrightarrow\;
 A_g=A_f\ \land\ \Lambda_g=\Lambda_f\ \land\
 \forall x:\mathrm{Morphism}(S,T),\ \Cov(x,J_g)\leftrightarrow\Cov(x,J_f).
\end{equation}
Продолжение $c=(C,f,R;\pi)$ содержит доказательство $\pi$ корректности и границы потенциала $P(C,f)\le R$. Резерв $R$ --- достаточная верхняя граница ресурсов завершения, а не расход. Определим $b(c)=\bill(c)$ как фактическую модельную стоимость исчерпывающего исполнения текущего статического продолжения из остатка $R$ и нулевых дополнительных расходов. По финансируемому завершению $b(c)\le R$. Действующее состояние также требует $R\le B$, где $B$ --- доступный остаток средств.

\begin{theorem}[Замена с сохранением продолжения]\label{thm:rebase}
Замена $\mathcal R$ другой сертифицированной библиотекой $\mathcal R'$ при том же $\tau$, переинициализация ожидающих очередей правил и предоставление корректного нового резерва дают продолжение $c'$ с $f\simeq f'$. Его завершение имеет в точности то же каноническое членство, что и завершение $c$; членство обоих равно членству в $\CanonRef$.
\end{theorem}
\begin{proof}
Операция переинициализации сбрасывает только очередь правил ожидающего задания области портов. Она сохраняет его область, каждое задание действий и его необработанный суффикс, точно копирует $A,\Lambda$. Покрытие не зависит от очереди правил, поэтому инвариантно покрытие каждого задания и, следовательно, их списка. Отсюда \eqref{eq:same-cont}, а \eqref{eq:frontier-correct} переносится на новый фронтир. Финансируемое завершение каждого продолжения даёт членство в $\CanonRef$ независимо от выбранной библиотеки правил. Это \rc{rebase\_same}, \rc{same\_correct} и \rc{completion\_full} в \rc{ControllerReplacement.lean}.
\end{proof}
Эта эквивалентность сильнее равенства множеств конечных ответов: она сохраняет денотацию оставшихся обязательств без перезапуска поиска. Она не разрешает произвольную миграцию исполнимого кода или перенос полного состояния обучения.

\subsection{Ликвидность и сравнительные свидетельства стоимости}
Абсолютное предложение индексировано точным текущим продолжением $c$ и содержит новый контроллер, доказательства неизменности тарифов и финансирования, плату за миграцию $K$ и доказанные границы $L\le b(c)$ и $b(c')\le U$. Пусть $H$ --- плата за проверку. Числовой механизм допуска принимает предложение тогда и только тогда, когда
\begin{equation}\label{eq:absolute-gate}
 H+R\le B,\qquad H+K+R'\le B,\qquad H+K+U<L.
\end{equation}
Первое неравенство защищает завершение текущим контроллером даже при неудаче проверки. Второе защищает завершение после миграции. Последнее устанавливает строгий выигрыш с учётом платы за проверку и миграцию. Низкая оставшаяся стоимость не означает доступности замены, поскольку резерв отличен от стоимости.

Независимые границы могут терять существенные корреляции. Поэтому конструкция относительных свидетельств заменяет одну пару $(L,U)$ конечным списком $E$ совместных пар нижних и верхних границ с предоставленным доказательством покрытия действительной пары:
\begin{equation}\label{eq:actual-coverage}
 \exists (\ell,u)\in E:\quad \ell\le b(c)\quad\land\quad b(c')\le u.
\end{equation}
Это совместный сертификат, а не список несвязанных маргинальных границ. В частности, $E$ не может быть пустым на типизированной границе, хотя необработанный универсальный арифметический тест принял бы пустой список. Предикат стоимости имеет вид
\begin{equation}\label{eq:relative-gate}
 \operatorname{relativeOK}(H+K,E)\;\Longleftrightarrow\;
 \forall(\ell,u)\in E,\ H+K+u<\ell.
\end{equation}
Для одноэлементного списка $[(L,U)]$ восстанавливается абсолютный экономический тест \eqref{eq:absolute-gate}. Например, пары $(7,1)$ и $(70,64)$ проходят при $H+K=2$, тогда как независимые экстремумы $(7,64)$ не проходят. Вторая строка свидетеля релиза представляет синтетическую неопределённость, а не измеренный второй мир.

\begin{theorem}[Точный допуск и корректность выигрыша]\label{thm:gate}
Для типизированного относительного предложения фиксированный числовой механизм допуска возвращает \emph{accepted} (принято) тогда и только тогда, когда выполнены первые два неравенства \eqref{eq:absolute-gate} и \eqref{eq:relative-gate}. При принятии $H+K+b(c')<b(c)$. Применение сохраняет эквивалентность продолжения, базовые тарифы поиска и финансируемый резерв.
\end{theorem}
\begin{proof}
Эквивалентность следует из последовательных ветвей \rc{gate}; универсальный тест строк отражён в булево значение посредством \rc{relativeOK\_iff}. Выберем покрывающую пару из \eqref{eq:actual-coverage}. Тогда $H+K+b(c')\le H+K+u<\ell\le b(c)$, что доказывает строгий выигрыш. Переинициализация использует теорему~\ref{thm:rebase}. Остаток после принятия $B-H-K$ не меньше $R'$ по второму условию. При отклонении сохраняется прежнее продолжение, а первое условие ограничивает его резерв остатком после проверки. При отсрочке состояние не меняется. Это \rc{gate\_accepted\_iff}, \rc{relative\_sound}, \rc{apply\_same} и \rc{apply\_liquid}.
\end{proof}

Публичная попытка использует $H=h+1$ для базовой платы $h\in\Nat$. Её вердикт --- обычный четырёхзначный результат; применение требует доказательства, что этот вердикт является выходом фиксированного механизма допуска для данного точного состояния и предложения. Таблица~\ref{tab:debit} задаёт операционный учёт. Первое нарушенное условие определяет вердикт; тест строк вычисляется только после прохождения обоих условий резерва.
\begin{table}[htbp]
\centering\small
\caption{Учёт допуска. $d$ --- списание, а не резерв завершения.}\label{tab:debit}
\begin{tabularx}{\linewidth}{@{}lXll@{}}\toprule
Вердикт & Условие (по порядку) & Установлено & $d$\\\midrule
Отсрочка & Не выполнено условие прежнего резерва & $c$ & $0$\\
Отказ: ликвидность & Не выполнено условие нового резерва & $c$ & $H$\\
Отказ: выигрыш & Не пройден тест совместных пар & $c$ & $H$\\
Принято & Пройдены все три теста & $c'$ & $H+K$\\\bottomrule
\end{tabularx}
\end{table}

Пусть сессия $s$ содержит действующее состояние и относительный журнал. Обозначим $D(s)=\spent(s)$ как сумму поисковых квитанций, плат за прежние попытки и плат относительного журнала, а $W(s)=B(s)+D(s)$ как общие средства. Применение вычитает ровно $d$ из остатка и добавляет $d$ в журнал, сохраняя $W$. Предложение уже несёт доказательства; механизм допуска не оценивает неизвестные стоимости бесплатными пробными запусками и не ищет доказательств. Он проверяет объявленные числовые условия. Внешние платежи за свидетельства регистрируют расход, но сами по себе не удостоверяют истинность или приобретение свидетельства.

\subsection{Закон смешанных историй замен}
Для финансируемой сессии определим целевую функцию завершения
\begin{equation}\label{eq:root-objective}
 \Phi(s)=D(s)+b(c_s).
\end{equation}
Это полные модельные расходы при завершении текущего продолжения, а не онлайн-оценка. Разрешённые действия --- финансируемый префикс текущего поиска, попытка, индексированная текущим продолжением, и доступная по цене внешняя квитанция о свидетельстве. Связная трасса $s_0\to\cdots\to s_n$ состоит из фактических применений действий с равенством соседних концов; это не список произвольно заявленных постусловий. В частности, разрешимое равенство произвольных состояний, содержащих функции, не предполагается.

Для каждого шага определим чистый выигрыш $g_i$ и накладные расходы $o_i$. Принятая замена имеет $g_i=\Phi(s_{i-1})-\Phi(s_i)>0$ уже за вычетом проверки и миграции, и $o_i=0$. Оплаченный отказ или плата за свидетельство имеет $g_i=0$, а $o_i$ равно соответствующему платежу. У префикса поиска и отсрочки обе величины нулевые. Пусть $G=\sum_i g_i$, $O=\sum_i o_i$, $N_p$ --- число неотложенных попыток замены, $N_a$ --- число принятых.

\begin{theorem}[Композиционный учёт корневых замен]\label{thm:root-trace}
Для каждой связной конечной корневой трассы
\begin{align}
 \Phi(s_n)+G&=\Phi(s_0)+O,\label{eq:trace-objective}\\
 W(s_n)&=W(s_0),& R(s_n)&\le B(s_n),\label{eq:trace-funds}\\
 N_p+B(s_n)&\le B(s_0),& N_a&\le G.
\end{align}
Кроме того, отношение концов $\Pres(s_0,s_n)$ точно сохраняет базовые тарифы поиска и прежнюю историю попыток, список сертификатов --- как префикс, а поисковые квитанции и относительный журнал --- как суффиксы списков в обратном хронологическом порядке.
\end{theorem}
\begin{proof}
Для префикса фиксированного контроллера детерминированное продолжение даёт $b(c)=d+b(c')$, где $d$ --- затраты префикса. Поэтому \eqref{eq:root-objective} не меняется. Отложенная попытка тождественна. Отклонённая попытка или квитанция свидетельства сохраняет $c$ и увеличивает $D$ на записанный платёж. Принятие строго уменьшает $\Phi$ по теореме~\ref{thm:gate}; эта разность определяет чистый выигрыш. Следовательно, каждый шаг удовлетворяет $\Phi(s_i)+g_i=\Phi(s_{i-1})+o_i$, сохранению средств и указанному отношению истории. Отношения префикса сертификатов и суффикса квитанций компонуются конкатенацией списков. Индукция по связной трассе складывает обе неотрицательные статистики на каждом шаге; положительные тарифы проверки ограничивают число оплаченных попыток, каждый принятый выигрыш не меньше единицы. Тождества: \rc{remainingBill\_prefix}, \rc{action\_laws}, \rc{laws\_compose} и \rc{trace\_laws} в \rc{RecursiveTrace.lean}.
\end{proof}
Поскольку $D(\finish(s))=\Phi(s)$, \eqref{eq:trace-objective} даёт необходимое и достаточное условие строгой чистой экономии относительно немедленного завершения текущим контроллером:
\begin{equation}\label{eq:net-gain}
 D(\finish(s_n))<D(\finish(s_0))\quad\Longleftrightarrow\quad O<G.
\end{equation}
Это \rc{trace\_final\_gain\_iff}. Теорема допускает убытки от исследовательских попыток без ослабления корректности. Она выявляет следующее препятствие: сам выбор момента таких попыток может быть неадекватен. Это мотивирует явное представление политики предложений как заменяемого объекта вместо фиксированной последовательности эвристических вмешательств.
